\documentclass[12pt,a4paper]{article}
\usepackage[utf8]{vietnam}
\usepackage{amsmath,amssymb}
\usepackage{tabularx}
\usepackage{tikz}
\usepackage{xcolor}
\usepackage[top=0.6cm, bottom=0.6cm, left=0.6cm, right=0.6cm]{geometry}
\usepackage[dethi]{ex_test}

\SAOPTN{widthSA=2.6,heightSA=0.62}
\setlength{\parskip}{0pt}

% Dong ke cham mo chuan viet tay
\newcommand{\dongke}[1]{%
  \begingroup%
  \setlength{\parskip}{0pt}%
  \baselineskip=18.5pt%
  \vspace{1.0mm}%
  \foreach \i in {1,...,#1}{%
    \noindent{\color{gray!60}\dotfill}\par%
  }%
  \vspace{0.6mm}%
  \endgroup%
}

% Dinh nghia o bai tap tu luan trong bang
\newcommand{\oBaiTap}[3][6]{%
  \vspace{1.0mm}%
  \noindent\textbf{#2} #3\par%
  \dongke{#1}%
}

\begin{document}
\pagestyle{empty}
% =============================================================================
% ========================= MA DE 101 - MAT TRUOC (TRANG 1/2) =============
% =============================================================================
\noindent
\begin{tabularx}{\linewidth}{|X|c|}
\hline
\textbf{PHIẾU KIỂM TRA ĐỊNH KỲ TOÁN 8 - TUẦN 08 - BUỔI 01} & \textbf{Mã đề: 101 - Trang 1/2} \\
\hline
\end{tabularx}

\vspace{1.2mm}

\noindent
\begin{tabularx}{\linewidth}{|p{3.2cm}|X|}
\hline
\centering\textbf{Điểm số} & \centering\textbf{Lời nhận xét của thầy cô} \tabularnewline
\rule{0pt}{1.8cm} & \\
\hline
\end{tabularx}

\vspace{1.2mm}

\noindent\textbf{Câu 1 (0,25 điểm).} Bậc của đơn thức $A = 5x^3 y^2 z$ bằng bao nhiêu? \shortans{$6$}\par\vspace{0.3mm}
\noindent\textbf{Câu 2 (0,25 điểm).} Bậc của đa thức $P = 3x^2 y - 5xy + 2x^2 y + 5xy - 7$ sau khi thu gọn là: \shortans{$3$}\par\vspace{0.3mm}
\noindent\textbf{Câu 3 (0,25 điểm).} Giá trị của $M = 2x^2 y - 3xy$ tại $x = 1, y = -2$ là: \shortans{$2$}\par\vspace{0.3mm}
\noindent\textbf{Câu 4 (0,25 điểm).} Trong khai triển $(x + 4)^2 = x^2 + ax + 16$, hệ số $a$ bằng: \shortans{$8$}\par\vspace{0.3mm}
\noindent\textbf{Câu 5 (0,25 điểm).} Tính nhanh giá trị biểu thức $65^2 - 35^2$ bằng: \shortans{$3000$}\par\vspace{0.3mm}
\noindent\textbf{Câu 6 (0,25 điểm).} Phân tích đa thức $6x^2 y - 9xy^2$ ta được $3xy(2x - \ldots)$. Điền vào chỗ trống: \shortans{$3y$}\par\vspace{0.3mm}
\noindent\textbf{Câu 7 (0,25 điểm).} Kết quả phép chia $(15x^4 y^3) : (3x^2 y)$ là: \shortans{$5x^2 y^2$}\par\vspace{0.3mm}
\noindent\textbf{Câu 8 (0,25 điểm).} Tứ giác $ABCD$ có $\widehat{A}=70^\circ, \widehat{B}=110^\circ, \widehat{C}=80^\circ$. Số đo $\widehat{D}$ là: \shortans{$100^\circ$}\par\vspace{0.3mm}
\noindent\textbf{Câu 9 (0,25 điểm).} Cho hình thang cân $ABCD$ ($AB \parallel CD$) có $\widehat{C} = 65^\circ$. Số đo góc $\widehat{D}$ là: \shortans{$65^\circ$}\par\vspace{0.3mm}
\noindent\textbf{Câu 10 (0,25 điểm).} Mua $10$ quyển vở giá $12\,000$ đ/quyển và $4$ bút giá $5\,000$ đ/cây. Tổng tiền là: \shortans{$140\,000\text{ đ}$}\par\vspace{0.3mm}

\vspace{0.5mm}
\noindent\textbf{Câu 11 (2,0 điểm). Thực hiện phép tính (hợp lý nếu có thể):}
\vspace{0.5mm}

\noindent
\begin{tabularx}{\linewidth}{|X|X|}
\hline
\oBaiTap[10]{a)}{$3x(x^2 - 2x + 1) - x^2(3x - 5)$} & \oBaiTap[10]{b)}{$(x + 3)^2 - (x - 3)(x + 3)$} \\
\hline
\oBaiTap[10]{c)}{$(12x^4 y^3 - 8x^3 y^2 + 4x^2 y) : 4x^2 y$} & \oBaiTap[10]{d)}{$(x - 2)(x + 5) - x(x + 3)$} \\
\hline
\end{tabularx}

\newpage
% =============================================================================
% ========================= MA DE 101 - MAT SAU (TRANG 2/2) ==============
% =============================================================================
\noindent
\begin{tabularx}{\linewidth}{|X|}
\hline
\centering\textbf{TOÁN 8 - TUẦN 08 - BUỔI 01 | Mã đề: 101 - Trang 2/2} \\
\hline
\end{tabularx}

\vspace{1.2mm}

\noindent\textbf{Câu 12 (2,0 điểm). Tìm $x$, biết:}
\vspace{0.5mm}

\noindent
\begin{tabularx}{\linewidth}{|X|X|}
\hline
\oBaiTap[5]{a)}{$2x(x - 4) - 2x^2 = 16$} & \oBaiTap[5]{b)}{$3x(x - 5) + 2(x - 5) = 0$} \\
\hline
\oBaiTap[5]{c)}{$(x + 2)^2 - 25 = 0$} & \oBaiTap[5]{d)}{$(2x - 3)^2 = 49$} \\
\hline
\end{tabularx}

\vspace{1.2mm}

\noindent
\begin{tabularx}{\linewidth}{|X|}
\hline
\vspace{0.5mm}
\textbf{Câu 13 (2,5 điểm).} Cho tam giác $ABC$ vuông tại $A$ ($AB < AC$), đường trung tuyến $AM$. Gọi $D$ là điểm đối xứng với $A$ qua $M$.\par
\textbf{a) (1,25 điểm)} Chứng minh tứ giác $ABDC$ là hình chữ nhật.\par
\textbf{b) (1,25 điểm)} Gọi $E$ là điểm đối xứng với $M$ qua $AC$, $ME$ cắt $AC$ tại $H$. Chứng minh tứ giác $AMCE$ là hình thoi.\par
\vspace{0.4mm}
\noindent\textit{Bài giải:}\dongke{11}
\\ \hline
\end{tabularx}

\vspace{1.2mm}

\noindent
\begin{tabularx}{\linewidth}{|X|}
\hline
\vspace{0.5mm}
\textbf{Câu 14 (1,0 điểm).} Cho biểu thức $A = x^2 - 6x + 2030$. Tìm giá trị nhỏ nhất của biểu thức $A$.\par
\vspace{0.4mm}
\noindent\textit{Bài giải:}\dongke{8}
\\ \hline
\end{tabularx}

\newpage
% =============================================================================
% ========================= MA DE 102 - MAT TRUOC (TRANG 1/2) =============
% =============================================================================
\noindent
\begin{tabularx}{\linewidth}{|X|c|}
\hline
\textbf{PHIẾU KIỂM TRA ĐỊNH KỲ TOÁN 8 - TUẦN 08 - BUỔI 01} & \textbf{Mã đề: 102 - Trang 1/2} \\
\hline
\end{tabularx}

\vspace{1.2mm}

\noindent
\begin{tabularx}{\linewidth}{|p{3.2cm}|X|}
\hline
\centering\textbf{Điểm số} & \centering\textbf{Lời nhận xét của thầy cô} \tabularnewline
\rule{0pt}{1.8cm} & \\
\hline
\end{tabularx}

\vspace{1.2mm}

\noindent\textbf{Câu 1 (0,25 điểm).} Bậc của đơn thức $B = -4x^2 y^3 z^2$ bằng bao nhiêu? \shortans{$7$}\par\vspace{0.3mm}
\noindent\textbf{Câu 2 (0,25 điểm).} Bậc của đa thức $Q = 4x^3 y - 2xy^2 + x^3 y + 2xy^2 + 1$ sau khi thu gọn là: \shortans{$4$}\par\vspace{0.3mm}
\noindent\textbf{Câu 3 (0,25 điểm).} Giá trị của $N = 3x^2 y - 2xy$ tại $x = -1, y = 2$ là: \shortans{$10$}\par\vspace{0.3mm}
\noindent\textbf{Câu 4 (0,25 điểm).} Trong khai triển $(x - 5)^2 = x^2 - bx + 25$, hệ số $b$ bằng: \shortans{$10$}\par\vspace{0.3mm}
\noindent\textbf{Câu 5 (0,25 điểm).} Tính nhanh giá trị biểu thức $75^2 - 25^2$ bằng: \shortans{$5000$}\par\vspace{0.3mm}
\noindent\textbf{Câu 6 (0,25 điểm).} Phân tích đa thức $8x^2 y - 12xy^2$ ta được $4xy(2x - \ldots)$. Điền vào chỗ trống: \shortans{$3y$}\par\vspace{0.3mm}
\noindent\textbf{Câu 7 (0,25 điểm).} Kết quả phép chia $(18x^5 y^4) : (6x^3 y^2)$ là: \shortans{$3x^2 y^2$}\par\vspace{0.3mm}
\noindent\textbf{Câu 8 (0,25 điểm).} Tứ giác $ABCD$ có $\widehat{A}=60^\circ, \widehat{B}=120^\circ, \widehat{C}=75^\circ$. Số đo $\widehat{D}$ là: \shortans{$105^\circ$}\par\vspace{0.3mm}
\noindent\textbf{Câu 9 (0,25 điểm).} Cho hình thang cân $ABCD$ ($AB \parallel CD$) có $\widehat{C} = 70^\circ$. Số đo góc $\widehat{D}$ là: \shortans{$70^\circ$}\par\vspace{0.3mm}
\noindent\textbf{Câu 10 (0,25 điểm).} Mua $8$ quyển vở giá $12\,000$ đ/quyển và $6$ bút giá $5\,000$ đ/cây. Tổng tiền là: \shortans{$126\,000\text{ đ}$}\par\vspace{0.3mm}

\vspace{0.5mm}
\noindent\textbf{Câu 11 (2,0 điểm). Thực hiện phép tính (hợp lý nếu có thể):}
\vspace{0.5mm}

\noindent
\begin{tabularx}{\linewidth}{|X|X|}
\hline
\oBaiTap[10]{a)}{$2x(x^2 - 3x + 2) - x^2(2x - 4)$} & \oBaiTap[10]{b)}{$(x + 4)^2 - (x - 4)(x + 4)$} \\
\hline
\oBaiTap[10]{c)}{$(15x^4 y^3 - 10x^3 y^2 + 5x^2 y) : 5x^2 y$} & \oBaiTap[10]{d)}{$(x - 3)(x + 7) - x(x + 4)$} \\
\hline
\end{tabularx}

\newpage
% =============================================================================
% ========================= MA DE 102 - MAT SAU (TRANG 2/2) ==============
% =============================================================================
\noindent
\begin{tabularx}{\linewidth}{|X|}
\hline
\centering\textbf{TOÁN 8 - TUẦN 08 - BUỔI 01 | Mã đề: 102 - Trang 2/2} \\
\hline
\end{tabularx}

\vspace{1.2mm}

\noindent\textbf{Câu 12 (2,0 điểm). Tìm $x$, biết:}
\vspace{0.5mm}

\noindent
\begin{tabularx}{\linewidth}{|X|X|}
\hline
\oBaiTap[5]{a)}{$3x(x - 3) - 3x^2 = 18$} & \oBaiTap[5]{b)}{$2x(x - 4) + 3(x - 4) = 0$} \\
\hline
\oBaiTap[5]{c)}{$(x + 3)^2 - 16 = 0$} & \oBaiTap[5]{d)}{$(2x - 5)^2 = 81$} \\
\hline
\end{tabularx}

\vspace{1.2mm}

\noindent
\begin{tabularx}{\linewidth}{|X|}
\hline
\vspace{0.5mm}
\textbf{Câu 13 (2,5 điểm).} Cho tam giác $ABC$ vuông tại $A$ ($AB < AC$), đường trung tuyến $AM$. Gọi $K$ là điểm đối xứng với $A$ qua $M$.\par
\textbf{a) (1,25 điểm)} Chứng minh tứ giác $ABKC$ là hình chữ nhật.\par
\textbf{b) (1,25 điểm)} Gọi $F$ là điểm đối xứng với $M$ qua $AB$, $MF$ cắt $AB$ tại $I$. Chứng minh tứ giác $AMBF$ là hình thoi.\par
\vspace{0.4mm}
\noindent\textit{Bài giải:}\dongke{11}
\\ \hline
\end{tabularx}

\vspace{1.2mm}

\noindent
\begin{tabularx}{\linewidth}{|X|}
\hline
\vspace{0.5mm}
\textbf{Câu 14 (1,0 điểm).} Cho biểu thức $B = x^2 - 8x + 2040$. Tìm giá trị nhỏ nhất của biểu thức $B$.\par
\vspace{0.4mm}
\noindent\textit{Bài giải:}\dongke{8}
\\ \hline
\end{tabularx}

\newpage
% =============================================================================
% ========================= MA DE 103 - MAT TRUOC (TRANG 1/2) =============
% =============================================================================
\noindent
\begin{tabularx}{\linewidth}{|X|c|}
\hline
\textbf{PHIẾU KIỂM TRA ĐỊNH KỲ TOÁN 8 - TUẦN 08 - BUỔI 01} & \textbf{Mã đề: 103 - Trang 1/2} \\
\hline
\end{tabularx}

\vspace{1.2mm}

\noindent
\begin{tabularx}{\linewidth}{|p{3.2cm}|X|}
\hline
\centering\textbf{Điểm số} & \centering\textbf{Lời nhận xét của thầy cô} \tabularnewline
\rule{0pt}{1.8cm} & \\
\hline
\end{tabularx}

\vspace{1.2mm}

\noindent\textbf{Câu 1 (0,25 điểm).} Bậc của đơn thức $C = 6x^4 y z^3$ bằng bao nhiêu? \shortans{$8$}\par\vspace{0.3mm}
\noindent\textbf{Câu 2 (0,25 điểm).} Bậc của đa thức $R = 2x^2 y^2 - 3x^3 + 3x^2 y^2 + 3x^3 - 9$ sau khi thu gọn là: \shortans{$4$}\par\vspace{0.3mm}
\noindent\textbf{Câu 3 (0,25 điểm).} Giá trị của $P = 2xy^2 - 4xy$ tại $x = 2, y = -1$ là: \shortans{$12$}\par\vspace{0.3mm}
\noindent\textbf{Câu 4 (0,25 điểm).} Trong khai triển $(x + 6)^2 = x^2 + cx + 36$, hệ số $c$ bằng: \shortans{$12$}\par\vspace{0.3mm}
\noindent\textbf{Câu 5 (0,25 điểm).} Tính nhanh giá trị biểu thức $85^2 - 15^2$ bằng: \shortans{$7000$}\par\vspace{0.3mm}
\noindent\textbf{Câu 6 (0,25 điểm).} Phân tích đa thức $10x^2 y - 15xy^2$ ta được $5xy(2x - \ldots)$. Điền vào chỗ trống: \shortans{$3y$}\par\vspace{0.3mm}
\noindent\textbf{Câu 7 (0,25 điểm).} Kết quả phép chia $(24x^4 y^5) : (4x^2 y^3)$ là: \shortans{$6x^2 y^2$}\par\vspace{0.3mm}
\noindent\textbf{Câu 8 (0,25 điểm).} Tứ giác $ABCD$ có $\widehat{A}=85^\circ, \widehat{B}=95^\circ, \widehat{C}=70^\circ$. Số đo $\widehat{D}$ là: \shortans{$110^\circ$}\par\vspace{0.3mm}
\noindent\textbf{Câu 9 (0,25 điểm).} Cho hình thang cân $ABCD$ ($AB \parallel CD$) có $\widehat{C} = 75^\circ$. Số đo góc $\widehat{D}$ là: \shortans{$75^\circ$}\par\vspace{0.3mm}
\noindent\textbf{Câu 10 (0,25 điểm).} Mua $12$ quyển vở giá $12\,000$ đ/quyển và $5$ bút giá $5\,000$ đ/cây. Tổng tiền là: \shortans{$169\,000\text{ đ}$}\par\vspace{0.3mm}

\vspace{0.5mm}
\noindent\textbf{Câu 11 (2,0 điểm). Thực hiện phép tính (hợp lý nếu có thể):}
\vspace{0.5mm}

\noindent
\begin{tabularx}{\linewidth}{|X|X|}
\hline
\oBaiTap[10]{a)}{$4x(x^2 - 2x + 1) - x^2(4x - 6)$} & \oBaiTap[10]{b)}{$(x + 5)^2 - (x - 5)(x + 5)$} \\
\hline
\oBaiTap[10]{c)}{$(18x^4 y^3 - 12x^3 y^2 + 6x^2 y) : 6x^2 y$} & \oBaiTap[10]{d)}{$(x - 4)(x + 6) - x(x + 2)$} \\
\hline
\end{tabularx}

\newpage
% =============================================================================
% ========================= MA DE 103 - MAT SAU (TRANG 2/2) ==============
% =============================================================================
\noindent
\begin{tabularx}{\linewidth}{|X|}
\hline
\centering\textbf{TOÁN 8 - TUẦN 08 - BUỔI 01 | Mã đề: 103 - Trang 2/2} \\
\hline
\end{tabularx}

\vspace{1.2mm}

\noindent\textbf{Câu 12 (2,0 điểm). Tìm $x$, biết:}
\vspace{0.5mm}

\noindent
\begin{tabularx}{\linewidth}{|X|X|}
\hline
\oBaiTap[5]{a)}{$4x(x - 2) - 4x^2 = 16$} & \oBaiTap[5]{b)}{$4x(x - 3) + 3(x - 3) = 0$} \\
\hline
\oBaiTap[5]{c)}{$(x + 4)^2 - 36 = 0$} & \oBaiTap[5]{d)}{$(2x - 1)^2 = 25$} \\
\hline
\end{tabularx}

\vspace{1.2mm}

\noindent
\begin{tabularx}{\linewidth}{|X|}
\hline
\vspace{0.5mm}
\textbf{Câu 13 (2,5 điểm).} Cho tam giác $MNP$ vuông tại $M$ ($MN < MP$), đường trung tuyến $MI$. Gọi $Q$ là điểm đối xứng với $M$ qua $I$.\par
\textbf{a) (1,25 điểm)} Chứng minh tứ giác $MNQP$ là hình chữ nhật.\par
\textbf{b) (1,25 điểm)} Gọi $K$ là điểm đối xứng với $I$ qua $MP$, $IK$ cắt $MP$ tại $H$. Chứng minh tứ giác $MIPK$ là hình thoi.\par
\vspace{0.4mm}
\noindent\textit{Bài giải:}\dongke{11}
\\ \hline
\end{tabularx}

\vspace{1.2mm}

\noindent
\begin{tabularx}{\linewidth}{|X|}
\hline
\vspace{0.5mm}
\textbf{Câu 14 (1,0 điểm).} Cho biểu thức $C = x^2 - 10x + 2050$. Tìm giá trị nhỏ nhất của biểu thức $C$.\par
\vspace{0.4mm}
\noindent\textit{Bài giải:}\dongke{8}
\\ \hline
\end{tabularx}

\newpage
% =============================================================================
% ========================= MA DE 104 - MAT TRUOC (TRANG 1/2) =============
% =============================================================================
\noindent
\begin{tabularx}{\linewidth}{|X|c|}
\hline
\textbf{PHIẾU KIỂM TRA ĐỊNH KỲ TOÁN 8 - TUẦN 08 - BUỔI 01} & \textbf{Mã đề: 104 - Trang 1/2} \\
\hline
\end{tabularx}

\vspace{1.2mm}

\noindent
\begin{tabularx}{\linewidth}{|p{3.2cm}|X|}
\hline
\centering\textbf{Điểm số} & \centering\textbf{Lời nhận xét của thầy cô} \tabularnewline
\rule{0pt}{1.8cm} & \\
\hline
\end{tabularx}

\vspace{1.2mm}

\noindent\textbf{Câu 1 (0,25 điểm).} Bậc của đơn thức $D = -7x^2 y^2 z$ bằng bao nhiêu? \shortans{$5$}\par\vspace{0.3mm}
\noindent\textbf{Câu 2 (0,25 điểm).} Bậc của đa thức $S = 5xy^2 - 4x^2 + 2xy^2 + 4x^2 - 3$ sau khi thu gọn là: \shortans{$3$}\par\vspace{0.3mm}
\noindent\textbf{Câu 3 (0,25 điểm).} Giá trị của $Q = x^2 y - 5xy$ tại $x = -2, y = 1$ là: \shortans{$14$}\par\vspace{0.3mm}
\noindent\textbf{Câu 4 (0,25 điểm).} Trong khai triển $(x - 7)^2 = x^2 - dx + 49$, hệ số $d$ bằng: \shortans{$14$}\par\vspace{0.3mm}
\noindent\textbf{Câu 5 (0,25 điểm).} Tính nhanh giá trị biểu thức $55^2 - 45^2$ bằng: \shortans{$1000$}\par\vspace{0.3mm}
\noindent\textbf{Câu 6 (0,25 điểm).} Phân tích đa thức $14x^2 y - 21xy^2$ ta được $7xy(2x - \ldots)$. Điền vào chỗ trống: \shortans{$3y$}\par\vspace{0.3mm}
\noindent\textbf{Câu 7 (0,25 điểm).} Kết quả phép chia $(28x^6 y^3) : (7x^4 y)$ là: \shortans{$4x^2 y^2$}\par\vspace{0.3mm}
\noindent\textbf{Câu 8 (0,25 điểm).} Tứ giác $ABCD$ có $\widehat{A}=65^\circ, \widehat{B}=115^\circ, \widehat{C}=85^\circ$. Số đo $\widehat{D}$ là: \shortans{$95^\circ$}\par\vspace{0.3mm}
\noindent\textbf{Câu 9 (0,25 điểm).} Cho hình thang cân $ABCD$ ($AB \parallel CD$) có $\widehat{C} = 80^\circ$. Số đo góc $\widehat{D}$ là: \shortans{$80^\circ$}\par\vspace{0.3mm}
\noindent\textbf{Câu 10 (0,25 điểm).} Mua $9$ quyển vở giá $12\,000$ đ/quyển và $8$ bút giá $5\,000$ đ/cây. Tổng tiền là: \shortans{$148\,000\text{ đ}$}\par\vspace{0.3mm}

\vspace{0.5mm}
\noindent\textbf{Câu 11 (2,0 điểm). Thực hiện phép tính (hợp lý nếu có thể):}
\vspace{0.5mm}

\noindent
\begin{tabularx}{\linewidth}{|X|X|}
\hline
\oBaiTap[10]{a)}{$3x(x^2 - 3x + 2) - x^2(3x - 7)$} & \oBaiTap[10]{b)}{$(x + 2)^2 - (x - 2)(x + 2)$} \\
\hline
\oBaiTap[10]{c)}{$(21x^4 y^3 - 14x^3 y^2 + 7x^2 y) : 7x^2 y$} & \oBaiTap[10]{d)}{$(x - 5)(x + 8) - x(x + 3)$} \\
\hline
\end{tabularx}

\newpage
% =============================================================================
% ========================= MA DE 104 - MAT SAU (TRANG 2/2) ==============
% =============================================================================
\noindent
\begin{tabularx}{\linewidth}{|X|}
\hline
\centering\textbf{TOÁN 8 - TUẦN 08 - BUỔI 01 | Mã đề: 104 - Trang 2/2} \\
\hline
\end{tabularx}

\vspace{1.2mm}

\noindent\textbf{Câu 12 (2,0 điểm). Tìm $x$, biết:}
\vspace{0.5mm}

\noindent
\begin{tabularx}{\linewidth}{|X|X|}
\hline
\oBaiTap[5]{a)}{$5x(x - 2) - 5x^2 = 20$} & \oBaiTap[5]{b)}{$5x(x - 2) + 2(x - 2) = 0$} \\
\hline
\oBaiTap[5]{c)}{$(x + 1)^2 - 49 = 0$} & \oBaiTap[5]{d)}{$(2x - 7)^2 = 9$} \\
\hline
\end{tabularx}

\vspace{1.2mm}

\noindent
\begin{tabularx}{\linewidth}{|X|}
\hline
\vspace{0.5mm}
\textbf{Câu 13 (2,5 điểm).} Cho tam giác $DEF$ vuông tại $D$ ($DE < DF$), đường trung tuyến $DM$. Gọi $K$ là điểm đối xứng với $D$ qua $M$.\par
\textbf{a) (1,25 điểm)} Chứng minh tứ giác $DEKF$ là hình chữ nhật.\par
\textbf{b) (1,25 điểm)} Gọi $P$ là điểm đối xứng với $M$ qua $DE$, $MP$ cắt $DE$ tại $I$. Chứng minh tứ giác $DMEP$ là hình thoi.\par
\vspace{0.4mm}
\noindent\textit{Bài giải:}\dongke{11}
\\ \hline
\end{tabularx}

\vspace{1.2mm}

\noindent
\begin{tabularx}{\linewidth}{|X|}
\hline
\vspace{0.5mm}
\textbf{Câu 14 (1,0 điểm).} Cho biểu thức $D = x^2 - 4x + 2030$. Tìm giá trị nhỏ nhất của biểu thức $D$.\par
\vspace{0.4mm}
\noindent\textit{Bài giải:}\dongke{8}
\\ \hline
\end{tabularx}

\newpage
\end{document}
